
\subsubsection{The Old Church Order and the Organization of the Congregation}

In the state church, the congregation has, in truth, disappeared. The state church’s royal clergy have official districts, with local and geographical boundaries, within which everyone belongs to the parish; all who are born within these boundaries are baptized; all who are baptized are confirmed, and so on.

This is sometimes, on solemn occasions, called “congregation”; most often, however, it goes under the modest name “parish,” which is again divided into main parishes and annexes.

In large parts of our people in America, this is the model that is regarded as exemplary for the organization of the congregation. And it is this form of congregational structure, fashioned in imitation of what is worst in the state church, that now presents itself under the high-sounding name: the old church order.

To a free-thinking observer it seems strange that an entire church body should choose this corruption of the congregation as both its goal and its means. At times one would like to believe that the situation is not as bad as the free-church people claim. But when the United Church, both through its chairman and through its annual meeting, declares and resolves that this is precisely what it wants, then one has no choice but to believe them.

And perhaps it is not so strange, when viewed from a worldly standpoint. The clergy are, according to the old church order, the real church; for their sake the congregation exists; the congregation is to support them, and they in turn are to be the shining lights in the darkness of the age, since they “sit on Moses’ seat.”

And if the clergy themselves are thoroughly permeated by the spirit and nature of the world, then it is altogether reasonable that they should be very well pleased with a purely worldly and external form of congregational organization.

Just as the old church order holds that the congregation exists to support the priest, so it also holds that the priest exists to conduct services and perform ministerial functions for the congregation. In religious matters the priest alone is active; the congregation is essentially receptive. Therefore the old church order is not only fitting for worldly priests, who desire the income and honor of the clerical office, but that same old church order is also entirely fitting for worldly congregations, who gladly wish to have the Sunday finery of religion laid over the impurity and worldliness of everyday life.

Thus, when we view the matter from this side, we do not wonder that “my people love to have it so.”

But that a church body should decide that this is the path on which it will walk—that is not astonishing; it is terrifying.

The few believing and spiritually enlightened men and women among our people will understand this: that when the old church order is thus approved, recommended, and practiced by hundreds of priests and congregations, it becomes a power that hinders revival and dulls spiritual life—an influence that cannot be overestimated.

We therefore believe that the United Church has been unfortunate in its choice of watchword. By dealing in "old wine" itself, and by urging others to persist in it, it has helped many to continue in unrighteousness; and the congregation has suffered harm, because it did not receive better counsel.

\subsubsection{A Turn in the Discussion}

It has been left to J. A. S.\footnote{J. A. S. is the author’s mark more frequently used by Professor Dr. Theol. J. A. Schmidt. — Ed.} to defend the old church order on behalf of the United Church. He is usually the one who must do what others prefer not to do. He has therefore also, to the best of his ability, endeavored to mount a defense; but it did not seem to fall to everyone’s taste, what he brought forward.

“Folkelbladet” has pointed out that it is remarkable that, in the matter of the congregation, alongside the Word of God, a new or other authority has been set up, namely “the old church order,” to which the United Church wishes to adhere. For this can only be explained in one of two ways: either there is no guidance in the Word of God concerning the congregation, or else the United Church is not willing to follow the guidance that is there.

At this point the discussion had put the old church order in a difficult position; for if it were the same as the congregation of the New Testament, then the Word of God would indeed be sufficient, and this assurance entirely superfluous. But if it were something other than the congregation of the New Testament, then this was highly unfortunate for the church order and for its defenders.

It now appears that at least Pastor Kildahl in Chicago is not satisfied with the situation that thus emerged. He therefore comes forward in the latest issue of “The Lutheran” with a new defense of the old church order; and he now asserts that the old church order is nothing other than the constitution of the United Church. And then it is no great challenge to prove that the constitution of the United Church, with its one-sided rule by majority, is altogether excellent.

Pastor Kildahl has admittedly remained silent for quite some time before coming forward with this information. And some time was indeed required before such a claim could even be believed among the people of the majority. For even if Pastor Kildahl himself believes the claim, that does not make it true.

If Chairman Høyme had meant the constitution of the United Church and the form of church governance determined therein, then he would surely have had the ability to say so himself.

If the annual meeting of the United Church last year had meant the constitution, then this meeting too would surely have had sufficient ability to state that the United Church continued to uphold its constitution.

If “The Lutheran” last year had known that the old church order was nothing other than the constitution of the body, then “The Lutheran” would have said so last year.

If Dr. Schmidt had arrived at the thought that the old church order was the constitution, then he would not have misspoken, whether about “old wine” or about many other matters.

The matter is this: Chairman Høyme speaks of the old church order in connection with the question of living or dead congregations, and in connection with the question of the preaching of the Word and the administration of the sacraments in the congregation, and in connection with the question of awakening or sleep in the congregations.

The matter is this: the Committee also understood the Chairman’s statement, on the basis of his report, as standing in the closest connection with the question of the congregation’s life and its sources. The same understanding was shared by the few men who, at last year’s annual meeting, objected to the expression “the old church order.” In this same way Dr. Schmidt and *The Lutheran* have expressed themselves the entire time, so that one can see that they are thinking of congregational life and congregational formation—the so-called ecclesiastical custom in the strict sense—when there is talk of “the old church order.”

Therefore, if Pastor Kildahl carefully reads page 18 of the United Church’s Annual Report for 1897, he will rightly see what is at issue here.

It is, however, a good sign that Pastor Kildahl does not defend the old church order in the sense in which Chairman Høyme, the annual meeting, and J. A. S. speak of it.

For whether one stands on one side or the other in the outward party divisions, this is certain: that all Christian people, and especially all Christian pastors in our congregations, feel keenly that “the old church order” is a terrible hindrance to spiritual life.

It is presumably out of this very feeling that Pastor Kildahl makes such great effort to show that he and his are the truly free‑church people.

We have always believed that Pastor Kildahl, despite all his curious newspaper articles, was in his heart on the right side in the main matter. We rejoice over everyone who in truth becomes free and free‑church.

\subsubsection{Does the Proclamation of the Gospel Suffer}

For Lutheran free Christians, the greatest of all ecclesial concerns is that the Gospel be proclaimed purely and rightly. Upon this depend faith and life. What is required in the Lutheran Church is a complete and unabridged proclamation of the Word of God, so that Law and Gospel are preached in their full extent and each in its proper place, so that they may accomplish their own work.

In so thoroughly monarchical a state church as, for example, the Norwegian, where the congregation has absolutely nothing to say, the influence of the “old church order” upon the proclamation of the Word itself is in fact rather small. A serious-minded pastor can, unhindered by the congregation to which and over which the king has appointed him, preach both Law and Gospel, both repentance and faith. Many God-fearing pastors do so; and even if they are hated for it, they are not driven away on that account. The strong arm of the law protects them.

It is otherwise with church discipline and the administration of the sacraments, two matters which in many ways belong closely together. In the state church there is no room for church discipline in the general sense; those who are members of the parish have the legal right to be members thereof, and they can demand access to make free use of the church’s goods. The believing pastor must submit to this and endure pangs of conscience unavoidably, until he is dulled.

This changes in part when the old church order is maintained in congregations that are independent of the state. What here is most likely to suffer first and most severely is undoubtedly the proclamation. One forms—after the manner of the state church—congregations of worldly people. A pastor is chosen and is to preach for and within these congregations. The members of the congregation are in theory Christians, since they are in the congregation; in practice they are worldly, since they do not have life in faith in the Son of God. The pastor must address them as Christians, treat them as Christians, preach to them as though they were Christians. And when it comes to the lowest point, he must preach in such a way as pleases and satisfies the worldly members of the congregation. If he does not please them, they do not pay.

This is the greatest and most serious danger in a continuation within the free church of the state church’s manner of building congregations. It entails danger for the proclamation of the Word itself. And if the edge of the sword is dulled, there is not much benefit in its use.

Did not the doctrine of the justification of the world arise precisely from the influence of worldly congregations on the proclamation?

But even if this sorrowful phenomenon is explained otherwise, daily experience teaches believers throughout the congregations that there are very few places where one does not perceive the influence of the worldly majority upon the proclamation, the administration of the sacraments, and church discipline. We will not mention all the sins which pastors here and there bear upon themselves through their funeral sermons and in part other occasional addresses. But even the proper proclamation of the Word itself is dulled and weakened to such a degree that it is generally understood that the unconverted and unbelieving are a group of people who preferably are far away from “our congregation.”

All earnest ministers will have felt the temptation that accompanies standing as a minister of the Word in a sluggish and worldly congregation. Many who began well have perhaps succumbed to that temptation.

Is it not right to ask: Is it quite certain that “the old church order” is an unqualified good?

In Norway there was once, at a church synod, discussion of introducing some change in the old church order by abolishing the “assistants” and establishing a “congregational council” in their place. Then a minister of the old school rose and said: “I believe we do best to hold to Scripture: ‘Hold fast to what you have, so that no one may take your crown!’” But an old advocate of freedom—he was a provost as well—rose and replied: “That is indeed good, to hold fast to what one has, when one has something; but when one has nothing, then it would surely be best to try to get hold of something.”

It seems much the same with “the old church order.” It would be best first to be certain that it was an unqualified good, before one was so eager for its preservation.


\subsubsection{Does the Christian Life Also Suffer?}

This is, we hope, the last time for a while that we shall need to return to “the old church order.”

We have seen at least something of its weaknesses—weaknesses that could be remedied by “having a congregation.” By this it is not said that there is not also much in “the old church order” that may be valuable even for “the congregation.” We are thinking of such things as the church year, the fixed preaching texts, various liturgical formularies, and the like.

What we wish to touch upon briefly, in conclusion, is the question whether Christianity suffers under the old church order. This time we do not mean the Christian religion in general, nor the Christian life within the old congregational structure with its geographical boundaries; rather, this time we mean the personal Christianity of the true Christians themselves. Does the Christianity of Christians, of believers, suffer as a consequence of the old church order?

It is often seen that a noble plant fails to thrive when it is set in poor soil or deprived of sunshine, light, and air. It does not bear full fruit.

In the same way, the Christianity of all believers is not always equally developed, equally pure and sure and clear. The conditions under which it grows may be too harsh and impoverished. The Savior speaks of such good and noble seed that falls among thorns.

There are many Christian people also within the enclosure of the old church order. But their Christianity suffers so much that it becomes weak and often even stunted.

There are many things that contribute to this. We can mention only the most important.

Since the whole church, according to the old-fashioned state-church conception, is essentially an arrangement or an institution in which certain outward means are placed in the hands of certain outward persons, so that these so-called officials are to bring the means to the members of the church, one will readily understand that within this church system there cannot be much talk of Spirit.

The pastor preaches the Word and administers the sacraments. The members of the congregation hear the Word and receive the Lord’s Supper. But of the work of the Holy Spirit in heart and life and labor, little or nothing is heard, because one does not want to have a congregation.

The true believers under this church order will therefore unavoidably come to suffer in that their Christianity is, as it were, arrested in its growth when it reaches a certain point. They have access to hear the Word and to believe the Word, to receive the Lord’s Supper with thanksgiving and joy in their hearts. And there it stagnates. It does not suit the church’s task as a “salvation institution” that it should go further. For the next step would lead into the power and anointing of the Holy Spirit—and what would the members of the congregation do with that, when they are not working but paying?

Those, therefore, who are not encompassed by the Christianity of the old church order are precisely those of whom Christ speaks, for example, in John 7:38: “Whoever believes … in me, out of his life there shall flow, as the Scripture says, living streams of water.” “Rivers of water,” it says in the Greek; and one may well ask whether there is much of that where the old church order is allowed to maintain itself in all strictness.

We very much doubt it. The preaching does not encourage it. Attention is not directed toward it. If signs begin to appear of the “living streams of water,” such things are readily checked by the means at the priest’s disposal: this so easily turns into fanaticism; it causes such disturbance in the congregation, and so forth. And very often the truly spiritual Christians have stood before the choice either of bringing all testimony and labor to an end, or of “going out” from the worldly organization of the congregation.

Where, therefore, “the old church order” among us has been strictly carried through in all its worldly and carnal lifelessness, there the result has also many times been this, that Methodists or Baptists or Mission Friends or “Free–free” have harvested the fruit of every spiritual stirring within the so‑called congregation.

Spirit is dangerous, perilous, uncomfortable for “old church order.” But Spirit creates the congregation and thrives in it.

Here is where the contrast enters. And this contrast is not from today or yesterday. This struggle has lasted long and has far yet to go. But certain it is that this is the first time in the history of the Christian church that an entire church body, by public decision, has crossed over to the side of mere form. And we cannot wonder that there surges through our people a storm which daily raises ever stronger protest, in awakened souls, against this system.

We may stop for this time. We know that there are many who, with us, are willing both to labor and to suffer for the restoration of the congregation in apostolic Spirit and truth.

Our hope is that the starling will rise, and that it will succeed at least for a portion of our Christian people to press forward toward the goal which, since Hans Nielsen Hauge, for more than a hundred years has stood admonishing the heart and mind of Norwegian Christians:

Do you want church?

\subsubsection{Congregation or Old Church Order}

It is a well-tested truth in human history that progress cannot occur without conflict, and that conflicts, as a rule, have their acute crises, in which the new that is pressing forward comes into a death struggle with the old that seeks to suffocate it. Then everything looks grim indeed. Is life to die, and death to prevail? Is the grave to close itself, black and dark, over the young, budding hope?

In not a few places among our people there are at present quite serious conflicts, in part violent crises. Shall the congregation, which has sought to break through, once again be suffocated by worldliness; or shall it be granted its springtime and its flourishing among us?

There are especially two enemies of the congregation that have appeared in our short church history. And as often as occasion has arisen, these two enemies have entered into alliance and joined forces with one another. They are clerical domination and worldliness. They are, after all, two shoots from the same root and stock. Together they constitute what, from certain quarters, has been called “the old church order.”

By “clerical domination” we mean here the abuse which a great theologian has described in this way: “Clerical domination seeks to turn the office of pastoral care into an unspiritual lordship over the congregation, and it does not shepherd the flock, but rejoices in its immaturity and seeks to exploit it for the sinful aims of worldly advancement and self-interest.”

By “worldliness” in the congregation we mean that disposition which seeks to make the congregation serve as a religious covering over a fleshly and worldly life, so that it may function to shield heart and conscience from the sharp goad of the truth, and even serve to guarantee eternal salvation to those who, with an unrepentant mind, have stood in outward connection with the congregation.

The kinship and friendship between clerical domination and worldliness are easy enough to see and recognize. The worldly disposition is quite willing to grant the priest both precedence and advantage, provided he spares the fleshly heart and does not turn the goad of the truth against it. And the fleshly priest is quite willing to let worldliness flourish in the congregation, provided it yields him the honor and advantage that he seeks.

This unholy alliance, which, as has been said, has from a certain quarter been adorned with the name “the old church order,” is of course the congregation’s most dangerous enemy. In times of calm it is like the “Great Burial Mound,” something heavy, slimy, and impassable, which everywhere presses against light and progress. But if anyone seriously demands the reality of the congregation where the name of congregation exists, then the dragon hisses and the serpent strikes and spews its poison over the striving combatant, however insignificant he may be.

Yet the pastor is not always willing to bow to worldliness; nor is the congregation always willing to bow to clerical rule. Therefore the struggle assumes different forms and manifests itself in different ways. From this there has arisen among us the struggle over lay ministry, over the pulpit, and now most recently over the congregation itself. These are, after all, only different stages in the same development.

The battle positions are therefore also quite varied. There are pastors who manfully fight against clerical domination and worldliness, just as there are congregations that with visible satisfaction submit to both. And the reverse as well. But beneath all this the conflicts continue, and we are steadily moving forward toward a genuine congregational life, where the words of Jesus are taken seriously: “You are all brothers,” and “the greatest among you shall be the servant of all.”

Therefore this is no time to lose heart. It is moving forward; and the speed can be measured by the waves of resistance that crash against the boat. That signifies progress. Pull in the oars and drift with the current, and soon the waves will subside and the noise cease. That is the calm of death.

Yet the struggle is not always equally hard, even though it is an unceasing struggle, as long as life and death contend with one another. It is not only a matter of the congregation advancing; it is also a matter of its being preserved. The latter is no less than the former.

The victories that have thus far been won in our ecclesial work through the recognition of lay ministry and Christian pastoral formation are of subordinate importance if they are not followed by, and do not bear their fruit in, the liberation of congregational life from the yoke of secular domination.

Is there not here something to labor for? Ought not all Christians to take part in this work?

